\documentclass[10pt,a4paper]{amsart}
\usepackage[margin=19mm]{geometry}
\usepackage{amsmath,amssymb,mathtools}
\usepackage[scr=rsfs,cal=euler]{mathalfa}
\usepackage{xfrac}
\usepackage[unicode,hidelinks,bookmarks=false]{hyperref}
\newcommand{\ie}{i.e.\ }
\newcommand{\cf}{cf.\ }
\newcommand{\resp}{respectively\ }
\newcommand{\Subsection}{\subsection}
\providecommand{\nobreakdash}{\nobreak\hbox{-}}
\setlength{\emergencystretch}{2em}
\multlinegap0pt
\usepackage{amssymb}
\usepackage{enumerate}
\usepackage{algorithm}\usepackage{algpseudocode}
\usepackage{tikz}
\newtheorem{definition}{Definition}
\newtheorem{theorem}{Theorem}
\newtheorem{proposition}{Proposition}
\title{A note on the Goldberg-Thorp example in light of the classification of linear ill-posed problems in Banach spaces}
\author{Bernd Hofmann, Jens Flemming}
\date{Source: arXiv:2603.01183v1 (CC BY 4.0)}
\begin{document}
\maketitle

\noindent\emph{Source and licence.} A note on the Goldberg-Thorp example in light of the classification of linear ill-posed problems in Banach spaces. Version arXiv:2603.01183v1 (CC BY 4.0), distributed by arXiv under the Creative Commons Attribution 4.0 International licence, http://creativecommons.org/licenses/by/4.0/, as recorded in the arXiv licence metadata for this exact version and shown on the abstract page https://arxiv.org/abs/2603.01183v1. Full licence notice: \texttt{licenses/arxiv-2603.01183-CC-BY-4.0.txt}.\par\smallskip

\noindent\textit{Editorial scope.} Counted unit: the article's proposition eq:xmn (source0.tex lines 321--326). The article's complete original proof of it is delivered with it (source0.tex lines 328--357, one \texttt{proof} environment of 30 lines). No mathematics is written, repaired or re-derived: the statement and the proof are the article's own lines, in the article's own notation, and no retained line is substituted or re-flowed.  Retained uncounted as the source-local context the proof needs: lines 108--110; lines 114--116; lines 156--176; lines 272--288.  Nothing was omitted from any retained span, so the package carries no omission record.  No retained line contains a citation, so the wrapper carries no bibliography and no bibliography entry.  No retained line contains a figure environment, an image-inclusion command or drawing code, so no image is delivered and the package declares no asset dependency.  Editorial formatting only, in addition: an AMS \texttt{amsart} preamble that restates the article's own declarations and macros used by the retained lines, namely the environments \texttt{definition}, \texttt{theorem}, \texttt{proposition} on separate counters, together with the standard packages those lines need.  The retained environments print in this document as Definition 1, Theorem 1, Proposition 1 in order of appearance, whereas the article prints the same items with the numbering of its own full document.  The complete native source of the article as distributed in the arXiv e-print for this exact version is bundled unchanged as \texttt{sources/195546/source0.tex}.
\begin{equation}\label{eq:opeq}
A\, x =y  \quad (x \in X,\; y \in Y)\,,
\end{equation}

\begin{equation}\label{eq:Beq}
B\, x =y  \quad (x \in \ell^1,\; y \in \ell^2)
\end{equation}

\begin{definition}[Well- and ill-posedness characterization and classification] \label{def:new}
Let $A: X \to Y$ be a bounded linear operator mapping between the
infinite-dimensional Banach spaces $X$ and $Y$.

Then the operator equation \eqref{eq:opeq} is called \emph{well-posed} if
\begin{align*} &\text{the range $\mathcal{R}(A)$ of $A$ is a \emph{closed} subset of $Y$ and, moreover,} \\
&\text{the null-space $\mathcal{N}(A)$ is \emph{complemented} in $X$;}
\end{align*}
otherwise the equation (\ref{eq:opeq}) is called \emph{ill-posed}.

\medskip

In the ill-posed case, (\ref{eq:opeq}) is called \emph{ill-posed of type I}   if
\begin{align*}
&\text{the range $\mathcal{R}(A)$ contains an \emph{infinite-dimensional closed subspace} of $Y$;}
%\\&\text{the null-space $\mathcal{N}(A)$ is \emph{complemented} in $A^{-1}[M]$.}
\end{align*}
%Here, $A^{-1}[M]$ is  the inverse image of $M$ under $A$.
otherwise the ill-posed
equation (\ref{eq:opeq}) is called \emph{ill-posed of type~II}.
\end{definition}

\begin{theorem} \label{thm:Bprop}
Any surjective bounded linear operator $B: \ell^1 \to \ell^2$ is of hybrid-type, and hence the operator equation \eqref{eq:Beq} is ill-posed of type~I in the sense of Definition~\ref{def:new}.
The operator $B$ has a null-space $\mathcal{N}(B)$, which is uncomplemented in $\ell^1$. Both the dimension and the codimension of the null-space $\mathcal{N}(B)$ are not finite. Consequently,  $B$ is non-injective. The adjoint operator $B^*: \ell^2 \to \ell^\infty$ is injective and hence not strictly singular with an infinite-dimensional closed range $\mathcal{R}(B^*)=\mathcal{N}(B)^\perp$ (annihilator of $\mathcal{N}(B)$). Therefore, the adjoint operator equation
\begin{equation}\label{eq:Beq*}
B^*\,\eta  =z  \quad (\eta \in \ell^2,\; z \in \ell^\infty)
\end{equation}
is well-posed in the sense of Definition~\ref{def:new}.

If $B$ is of Mazur-type attaining the form
\begin{equation}\label{eq:Bstructure}
B\,x:=\sum_{k=1}^\infty x_k\,\zeta^{(k)} \in \ell^2 \qquad \mbox{for} \quad x=(x_1,x_2,...) \in \ell^1,
\end{equation}
the adjoint operator $B^*$ can be explicitly written as
\begin{equation}\label{eq:B*}
B^*\,\eta= (\langle\eta,\zeta^{(1)}\rangle_{\ell^2},\langle\eta,\zeta^{(2)}\rangle_{\ell^2},\ldots)   \in \ell^\infty.
\end{equation}
\end{theorem}

\begin{proposition}
If $y$ is not a multiple of $\zeta^{(k)}$ for some $k$, then there is no minimum-norm solution $x_{mn}$ satisfying
\begin{equation}\label{eq:xmn}
\|x_{mn}\|_{\ell^1} = \min_{\{x\in X:Bx=y\}}\|x\|_{\ell^1}.
\end{equation}
\end{proposition}

\begin{proof}
The minimization problem \eqref{eq:xmn} is equivalent to
\begin{equation*}
\mathcal{N}(B)^\perp\cap\partial\|\cdot\|_{\ell^1}(x_{mn})\neq\emptyset
\end{equation*}
with
\begin{equation*}
\xi\in\partial\|\cdot\|_{\ell^1}(x)\quad\Leftrightarrow\quad\xi_k\begin{cases}
=-1,&\text{if }x_k<0,\\
\in[-1,1],&\text{if }x_k=0,\\
=1,&\text{if }x_k>0
\end{cases}
\;\text{for }k\in\mathbb{N}.
\end{equation*}
From Theorem~\ref{thm:Bprop} we know $\mathcal{N}(B)^\perp=\mathcal{R}(B^\ast)$. Thus, if there exists a minimizer, then there is some $\eta\in\ell^2$ such that $B^\ast\eta\in\partial\|\cdot\|_{\ell^1}(x_{mn})$.
\par
Assume that $x_{mn}$ has at least two non-zero components $[x_{mn}]_m\neq 0$ and $[x_{mn}]_n\neq 0$. Denote the signs of both components by $s_m\in\{-1,1\}$ and $s_n\in\{-1,1\}$, respectively. Then $[B^\ast\eta]_m=s_m$ and $[B^\ast\eta]_n=s_n$ or, equivalently, $\langle\eta,\zeta^{(m)}_{\ell^2}\rangle=s_m$ and $\langle\eta,\zeta^{(n)}\rangle_{\ell^2}=s_n$. Now take a subsequence $(\zeta^{(k_l)})_{l\in\mathbb{N}}$ of $(\zeta^{(k)})_{k\in\mathbb{N}}$ converging to
\begin{equation}
\tilde{\eta}:=\frac{s_m\,\zeta^{(m)}+s_n\,\zeta^{(n)}}{\|s_m\,\zeta^{(m)}+s_n\,\zeta^{(n)}\|_{\ell^2}}.
\end{equation}
Then $\langle\eta,\zeta^{(k_l)}\rangle_{\ell^2}\to\langle\eta,\tilde{\eta}\rangle_{\ell^2}$. From
\begin{equation}
\langle\eta,\tilde{\eta}\rangle_{\ell^2}=\frac{2}{\|s_m\,\zeta^{(m)}+s_n\,\zeta^{(n)}\|_{\ell^2}}
\end{equation}
we see $\langle\eta,\tilde{\eta}\rangle_{\ell^2}\geq 1$ and that $\langle\eta,\tilde{\eta}\rangle_{\ell^2}>1$ holds if and only if $\zeta^{(m)}$ and $\zeta^{(n)}$ are linearly dependent. Thus, $\langle\eta,\tilde{\eta}\rangle_{\ell^2}=1$ is only possible for $\zeta^{(n)}=-\zeta^{(m)}$.
\par
In case $\langle\eta,\tilde{\eta}\rangle_{\ell^2}>1$ we find (large enough) $k_l$ with $\langle\eta,\zeta^{(k_l)}\rangle_{\ell^2}>1$ or, equivalently $[B^\ast\eta]_{k_l}~>~1$. Thus, $B^\ast\,\eta\notin\partial\|\cdot\|_{\ell^1}(x_{mn})$, which shows that $x_{mn}$ can have at most one non-zero component $[x_{mn}]_k$ if $\zeta^{(l)}\neq-\zeta^{(k)}$ for all $l\neq k$. With only one non-zero component in $x_{mn}$ we obtain $Bx_{mn}=[x_{mn}]_k\zeta^{(k)}$.
\par
In case $\langle\eta,\tilde{\eta}\rangle_{\ell^2}=1$, we do not obtain a contradiction (at the moment). That is, $x_{mn}$ may have two non-zero components $[x_{mn}]_m$ and $[x_{mn}]_n$ as long as $\zeta^{(n)}=-\zeta^{(m)}$. But a third non-zero component $[x_{mn}]_l$ is not possible, because corresponding $\zeta^{(l)}$ would have to be equal to both $-\zeta^{(m)}$ and $-\zeta^{(n)}=\zeta^{(m)}$. Thus, $Bx_{mn}=([x_{mn}]_m-[x_{mn}]_n)\zeta^{(m)}$.
\end{proof}

\end{document}
